%% ===================================================================
%% Appendix C. From Continuous Sensitivity to Discrete Rankings
%% Proofs of prop:transfer (continuous-to-discrete ranking transfer, with
%% the full second-order remainder) and prop:explicit (K-layer explicit
%% GNNs). Labels app:proof_transfer / eq:ranking / app:explicit_prop are
%% referenced from the main text and preserved.
%% ===================================================================
\section{From Continuous Sensitivity to Discrete Rankings}
\label{app:transfer}

The sensitivity $S_c$ is a continuous, implicit-model object, while practitioners delete whole edges and
use explicit, finite-depth GNNs. Two bridges carry $S_c$ across that gap, and we prove both here: the
first shows the continuous score ranks discrete edge removals, the second shows the same construction
applies to unrolled explicit networks.

\subsection{Continuous-to-Discrete Ranking Transfer}
\label{app:proof_transfer}

\Cref{prop:transfer} states that, when every single-edge removal is subcritical, the true removal damage
of edge $k$ factors as $d_k=w_k v_k+R_k$ with $\abs{R_k}\le L_J w_k^2/\big(2(1-\kappa)^2\big)$, so the
edge-weighted score $w_k v_k$ preserves the brute-force removal order whenever the score gap between two
edges beats this quadratic remainder. The pairwise condition under which a pair $(k_1,k_2)$ with
$v_{k_1}>v_{k_2}$ and $w_{k_1}\ge w_{k_2}$ keeps its order is
\begin{equation}
w_{k_1}<\frac{v_{k_1}-v_{k_2}}{L_J/(1-\kappa)^2}.
\label{eq:ranking}
\end{equation}

\begin{proof}
\emph{Step 1: a removal is a continuous edge-weight perturbation.}
Deleting edge $k$ under fixed normalization sets $[\delta\Ahat]_{ij}=[\delta\Ahat]_{ji}=-w_k$ with the
degree matrix $D$ held fixed, which is exactly the continuous perturbation $S_c$ is built to score. By
the first-order sensitivity \eqref{eq:ift-expand} restricted to the edge, the damage is
$\norm{\Delta\zstar}\approx w_k v_k$ with $v_k=\norm{[S_c]_{:,k}}_2$. (Recomputing the normalization
instead adds an $O(d_i^{-1})$ rescaling of the incident edges, negligible away from low-degree nodes.)
It remains to bound the second-order error $R_k$.

\emph{Step 2: the remainder is governed by one cross term, with nothing omitted.}
On the interior of a ReLU linear region the activation mask $D=\diag(\phi')$ is constant, so the operator
$F(z,\Ahat)=D\,(\Ahat\otimes W)\,z+Dx$ is \emph{bilinear} in $(\Ahat,z)$: linear in $z$ at fixed
$\Ahat$, and linear in $\Ahat$ at fixed $z$. Both pure second partials therefore vanish,
\begin{equation}
F_{zz}=\frac{\partial^2 F}{\partial z^2}=0,\qquad
F_{AA}=\frac{\partial^2 F}{\partial\mathrm{vec}(A)^2}=0,
\label{eq:bilinear}
\end{equation}
and the only surviving second derivative is the adjacency--state cross term $F_{zA}$. Differentiating
the IFT sensitivity $\partial\zstar/\partial\mathrm{vec}(A)=(I-J_z)^{-1}J_A$ once more and using
\eqref{eq:bilinear} to drop every other term gives the full Hessian of the equilibrium map,
\begin{equation}
\frac{\partial^2\zstar}{\partial\mathrm{vec}(A)^2}
=(I-J_z)^{-1}\big(F_{zA}+F_{Az}\big)(I-J_z)^{-1}J_A.
\label{eq:hessian-full}
\end{equation}
Each summand carries two resolvents $(I-J_z)^{-1}$, a cross partial $F_{zA}$ or $F_{Az}$ proportional to
$\norm{W}_2$, and $J_A$ proportional to $\norm{W}_2\norm{\zstar}$ (since
$J_A=\partial F/\partial\mathrm{vec}(A)\propto\zstar W^\top$). Bounding the resolvents by $1/(1-\kappa)$
each (\cref{eq:neumann-bound}) and collecting the weight factors into
\begin{equation}
L_J\le\norm{W}_2^2\,\norm{\zstar},\qquad \norm{\zstar}\le\frac{\norm{X_{\mathrm{proj}}}}{1-\kappa},
\label{eq:LJ-bound}
\end{equation}
where the bound on $\norm{\zstar}$ is the $\kappa$-contraction of $F$ at its fixed point (A3), yields
$\norm{\partial^2\zstar/\partial\mathrm{vec}(A)^2}_2\le(1-\kappa)^{-2}L_J$. The feasible segment
$\Ahat\to\Ahat+\delta\Ahat$ crosses finitely many regions (simultaneous boundary crossings are
non-generic, measure zero), and the conservative IFT of \citet{bolte2021conservative} applies on each,
so summing the second-order Taylor remainder across them preserves the bound:
\begin{equation}
\abs{R_k}\le\tfrac12\,(1-\kappa)^{-2}L_J\,w_k^2.
\label{eq:Rk-bound}
\end{equation}
This is the \emph{complete} second-order remainder: the bilinearity \eqref{eq:bilinear} leaves no other
second-derivative term to bound, so \eqref{eq:Rk-bound} is exact on each region rather than a
dominant-term reduction.

\emph{Step 3: the score preserves the ranking.}
With $d_k=w_k v_k+R_k$ from Steps 1--2 and $C:=L_J/\big(2(1-\kappa)^2\big)$, two edges differ by
\begin{equation}
d_{k_1}-d_{k_2}\ \ge\ \big(w_{k_1}v_{k_1}-w_{k_2}v_{k_2}\big)-C\big(w_{k_1}^2+w_{k_2}^2\big).
\label{eq:rank-gap}
\end{equation}
When $v_{k_1}>v_{k_2}$ and $w_{k_1}\ge w_{k_2}$ the first-order gap is positive, and condition
\eqref{eq:ranking} is exactly what makes it exceed the summed quadratic remainders, so the sign of
$d_{k_1}-d_{k_2}$ matches that of the score gap and the order is preserved. Because the signal scales as
$O(w_k)$ and the error only as $O(w_k^2)$, the ranking is reliable for the small-weight edges of sparse
graphs; empirically $\abs{R_k}$ runs $2$--$10\times$ below $L_J w_k^2$, and the edge-weighted score
reaches Kendall $\tau=+0.996$ against brute-force single-edge removal on full-graph Amazon Photo
(\cref{sec:explicit_extension}).
\end{proof}

\subsection{Structural Sensitivity for $K$-Layer Explicit GNNs}
\label{app:explicit_prop}

\Cref{prop:explicit} states that an explicit $K$-layer network
$Z_K=g_K\circ\cdots\circ g_1(X,A)$, which has no fixed point, still admits the unrolled sensitivity
$S_K$ of \cref{eq:unrolled}, and that $S_K$ replaces $S$ throughout the pipeline.

\begin{proof}
\emph{Step 1: unroll and apply the chain rule.}
Write $Z_K$ as the composition of the $K$ layers and differentiate $\mathrm{vec}(Z_K)$ with respect to
$\mathrm{vec}(A)$. Each layer $l$ depends on $A$ both directly, through
$J_A^{(l)}=\partial g_l/\partial\mathrm{vec}(A)$, and indirectly through its input $Z_{l-1}$; propagating
the direct dependence of layer $l$ forward through the later layers' input Jacobians
$J_z^{(k)}=\partial g_k/\partial Z_{k-1}$ and summing over $l$ gives
\begin{equation}
S_K=\frac{\partial\,\mathrm{vec}(Z_K)}{\partial\,\mathrm{vec}(A)}
=\sum_{l=1}^{K}\Big(\prod_{k=l+1}^{K}J_z^{(k)}\Big)J_A^{(l)},
\label{eq:unrolled-restate}
\end{equation}
which is \eqref{eq:unrolled}.

\emph{Step 2: the first-order shift bound.}
Submultiplicativity of the spectral norm applied term by term in \eqref{eq:unrolled-restate} gives
$\sigma_1(S_K)\le\sum_{l=1}^{K}\big(\prod_{k=l+1}^{K}\norm{J_z^{(k)}}_2\big)\norm{J_A^{(l)}}_2$, and a
first-order Taylor expansion of $Z_K$ in $A$ turns this into the shift bound
$\norm{\Delta Z_K}_F\le\sigma_1(S_K)\,\norm{\delta A}_F+O(\norm{\delta A}^2)$.

\emph{Step 3: the implicit limit.}
For a weight-tied model every $\norm{J_z^{(k)}}_2\le\kappa$, so the sum is geometric,
\begin{multline}
\sigma_1(S_K)\le\norm{J_A}_2\sum_{j=0}^{K-1}\kappa^{\,j}=\norm{J_A}_2\,\frac{1-\kappa^{K}}{1-\kappa}\\
\xrightarrow[K\to\infty]{}\ \frac{\norm{J_A}_2}{1-\kappa},
\label{eq:explicit-limit}
\end{multline}
recovering the implicit bound of \cref{app:proof_phase}: unrolling a finite network is the truncated
version of the geometric series that the implicit resolvent $(I-J_z)^{-1}$ sums in closed form, so the
two agree as depth grows. Because $S_K$ enters the constrained construction $S_c$ and the per-node
radius $r_v$ (\cref{prop:radius}) in exactly the same place as $S$, the whole pipeline carries over;
only the closed-form $\ecrit$, which needs the contraction (A3), stays restricted to the implicit
subclass. Empirically the explicit models match the implicit tightness ($0.99$--$1.02$ across the
suite, \cref{app:explicit}).
\end{proof}
